% !TeX root = ../main.tex

\section{Conclusion}\label{sec:conclusion}

This study has examined how crypto key opinion leaders understand their motivations, credibility, and responsibilities. We conceptualise KOLs as hybrid figures: educators, informal financial advisors, entrepreneurs, and community stewards, whose credibility emerges from motivational orientations, cross-platform and on-chain practices, and collectively enforced norms. Drawing on self-determination theory and interviews with 13 KOLs, we showed how autonomy, competence, and relatedness needs are negotiated alongside monetisation pressures, token holdings, and shared financial exposure with followers. Our hybrid human--LLM workflow for thematic analysis demonstrates the potential of large language models to broaden candidate theme coverage while underscoring the necessity of human oversight to mitigate hallucinations and hidden biases.

\textbf{Regulatory implications.}
These insights carry practical implications at a moment when influencer regulation around financial content is tightening globally, with jurisdictions such as China, the European Union, and Singapore introducing stricter oversight of crypto-related promotion \cite{ESMA2023MiCAR,MAS2024crypto}. Our findings suggest that effective governance will need to account for the blurred boundaries between education, promotion, and speculation that define crypto KOL work. Rather than treating KOLs simply as advertisers or financial advisors, regulators may benefit from recognising the hybrid nature of their roles and the community dynamics that sustain their influence. Designing mechanisms that support transparent incentives, encourage responsible signalling, and reduce asymmetries of financial exposure (such as the disclosure badges, accuracy indicators, and interoperable trust profiles discussed in Section~\ref{sec:discussion}) may strengthen both creator credibility and user protection as crypto markets continue to evolve.